\documentclass{article}
\usepackage[UTF8]{ctex}
\usepackage{amsmath,amssymb,amsfonts}
\usepackage{enumitem}
\usepackage{geometry}
\geometry{a4paper, margin=1in}
% === 错题本排版优化（不参与题目标识） ===
\setlength{\parindent}{0pt}
\setlength{\parskip}{0.35em plus 0.10em minus 0.08em}
\setlength{\emergencystretch}{3em}
\setlist{itemsep=0.25em, topsep=0.40em, parsep=0pt}
\allowdisplaybreaks[2]
% === 错题本排版优化结束 ===
\begin{document}
\section*{云南大学2025年数学分析真题}
\section*{194．云南大学2025年数学分析试卷}

1．极限 \(\lim _{x \rightarrow 3} \frac{\sqrt{x^{3}+9}-6}{2-\sqrt{x^{3}-23}}=\)

2．积分 \(\int_{2}^{+\infty} \frac{5}{x^{4}+3 x^{3}-4} \mathrm{~d} x=\) \(\_\_\_\_\)．

3．曲面 \(z=\frac{x^{2}}{2}+y^{2}-2\) 的平行于平面 \(2 x+2 y-z=0\) 的切平面方程是 \(\_\_\_\_\)．

4．设连续函数 \(f(x)\) 满足 \(f(x+2)-f(x)=x, \int_{0}^{2} f(x) \mathrm{d} x=0\) ，则 \(\int_{1}^{3} f(x) \mathrm{d} x=\) \(\_\_\_\_\)．

5．设 \(F(x)=\int_{0}^{x}\left(\int_{0}^{y^{2}} \frac{\sin t}{1+t^{2}} \mathrm{~d} t\right) \mathrm{d} y\) ，则 \(F^{\prime \prime}(x)=\) \(\_\_\_\_\)．

6．级数 \(\sum_{n=0}^{\infty} \frac{x^{2 n}}{(2 n)!}\) 的和函数是 \(\_\_\_\_\)．

二．设数列 \(\left\{x_{n}\right\}\) 满足：\(x_{1}>0, x_{n} e^{x_{n+1}}=e^{x_{n}}-1, n=1,2, \cdots\) ．

\begin{itemize}
  \item[（1）] 证明数列 \(\left\{x_{n}\right\}\) 收敛。
  \item[（2）] 求 \(\lim _{n \rightarrow \infty} x_{n}\) ．
\end{itemize}

三．若函数 \(f(x)\) 在 \([a, b]\) 上连续且单调增加．证明： \(\int_{a}^{b} x f(x) \mathrm{d} x \geq \frac{a+b}{2} \int_{a}^{b} f(x) \mathrm{d} x\) ．

\begin{itemize}
  \item[四．] 已知 \(f(x)=\int_{1}^{x} e^{u^{2}} \mathrm{~d} u\) ．

  \item[1．] 证明：存在 \(\xi \in(1,2)\) ，使得 \(f(\xi)=(2-\xi) e^{\xi^{2}}\) ．

  \item[2．] 证明：存在 \(\eta \in(1,2)\) ，使得 \(f(2)=\ln 2 \cdot \eta \cdot e^{\eta^{2}}\) ．
  \item[] 五。讨论当 \(p\) 取得何值时，级数 \(\sum_{n=1}^{\infty} \frac{(-1)^{n}}{n^{p+\frac{1}{n}}}\) 发散？条件收敛？绝对收敛？
\end{itemize}

六。设 \(f_{n}(x)=x n^{\alpha} e^{-n x}, n=1,2, \cdots\) 。

\begin{itemize}

  \item[1．] 当 \(\alpha\) 为何值时，函数序列 \(\left\{f_{n}(x)\right\}\) 在［0，1］上收敛？

  \item[2．] 当 \(\alpha\) 为何值时，函数序列 \(\left\{f_{n}(x)\right\}\) 在［0，1］上一致收敛？

  \item[3．] 当 \(\alpha\) 为何值时， \(\lim _{n \rightarrow \infty} \int_{0}^{1} f_{n}(x) \mathrm{d} x\) 和 \(\int_{0}^{1} \lim _{n \rightarrow \infty} f_{n}(x) \mathrm{d} x\) 不相等？
\end{itemize}

七．设

\[
f(x, y)= \begin{cases}x y \sin \frac{1}{x^{2}+y^{2}}, & x^{2}+y^{2} \neq 0 \\ 0, & x^{2}+y^{2}=0\end{cases}
\]

\begin{itemize}

  \item[1．] 证明 \(f(x, y)\) 的两个一阶偏导数在点 \((0,0)\) 不连续．

  \item[2．] 证明 \(f(x, y)\) 在 \((0,0)\) 可微．
\end{itemize}

八．计算三重积分

\[
I=\iiint_{\Omega} z \mathrm{~d} x \mathrm{~d} y \mathrm{~d} z
\]

其中积分区域 \(\Omega\) 是球面 \(x^{2}+y^{2}+z^{2}=8\) 的上半部分和抛物面 \(z=\frac{1}{2}\left(x^{2}+y^{2}\right)\) 所围立体区域．

九．设 \(L\) 是曲面 \(\sum: 4 x^{2}+y^{2}+z^{2}=1, x \geq 0, y \geq 0, z \geq 0\) 的边界，曲面方向朝上，曲线 \(L\) 的方向和曲面的方向符合右手法则，求曲线积分

\[
I=\oint_{L}\left(y z^{2}-\cos z\right) \mathrm{d} x+2 x z^{2} \mathrm{~d} y+(2 x y z+x \sin z) \mathrm{d} z .
\]
\end{document}
